\documentclass[11pt]{article}
\usepackage[T1]{fontenc}
\usepackage[utf8]{inputenc}
\usepackage{lmodern}
\usepackage[margin=1in]{geometry}
\usepackage{amsmath,amssymb,amsthm,mathtools}
\usepackage{enumitem}
\usepackage[colorlinks=true,linkcolor=black,citecolor=black,urlcolor=blue]{hyperref}
\usepackage{microtype}
\newtheorem{theorem}{Theorem}[section]
\newtheorem{proposition}[theorem]{Proposition}
\newtheorem{lemma}[theorem]{Lemma}
\newtheorem{corollary}[theorem]{Corollary}
\theoremstyle{definition}
\newtheorem{definition}[theorem]{Definition}
\newtheorem{example}[theorem]{Example}
\newtheorem{remark}[theorem]{Remark}
\newcommand{\Cat}{\mathbf{Cat}}
\newcommand{\Set}{\mathbf{Set}}
\newcommand{\Biprof}{\mathbf{Biprof}}
\newcommand{\Ps}{\operatorname{PsNat}}
\newcommand{\disc}{\operatorname{disc}}
\newcommand{\id}{\mathrm{id}}
\newcommand{\Ob}{\operatorname{Ob}}
\newcommand{\Kl}{\operatorname{Kl}}
\newcommand{\Ctx}{\operatorname{Ctx}}
\newcommand{\Cnt}{\operatorname{Cnt}}
\newcommand{\Prob}{\mathbf{Prob}^{+}_{\mathrm{fin}}}
\newcommand{\List}{\operatorname{List}}
\newcommand{\Psh}{\operatorname{Psh}}
\newcommand{\Copsh}{\operatorname{Copsh}}
\newcommand{\Tel}{\operatorname{Tel}}
\newcommand{\coend}[1]{\int_{\mathrm{ps}}^{#1}}
\newcommand{\Qrel}{\mathbb Q}
\newcommand{\then}{\mathbin{;}}
\title{One relative construction for context content and effects}
\author{}
\date{1 November 2026}
\begin{document}
\maketitle
\begin{abstract}
A single construction is given, rather than a list of related theories. Its
algebra object is a category-valued relative pseudopromonad in the existing
tricategory of biprofunctors. Its realization has the original computation
categories as hom categories and a designated interpretation of pure maps.
Companion bending recovers every changing-endpoint square. Context and content
presentations determine such objects; their relative tensor is proved globally
equivalent to the category of triples $(P,Q,eP\to dQ)$, using an explicit
pseudocoend normalization that does not invert grade maps. A supplied coherent
crossing makes this tensor one combined pseudopromonad, with a displayed unit
and multiplication. Dependent finite probability presentations and lists give
a fully worked nonthin instance with locally fully faithful one-sided maps.
Grothendieck extension and telescopes are internal or base-changing constructors
of the same system. Marked presentations recover the original fibres and
extensions; unmarked operational data are not claimed to recognize them
intrinsically. Crossings remain explicit compatibility data; a bicyclic obstruction rules
out universally conservative amalgamation.
The ambient, realization and composition machinery are established mathematics;
no historical novelty is claimed for that framework.
\end{abstract}

\section{The single structure and its scope}

The common structure is not the statement that everything is a bicategory.
It consists of a specified pure base $C$, one category-valued module $H$ over
that base, a relative unit and multiplication, and, when presentation-level
recovery is required, marked generic arrows. Composition is constructed in a
relative tensor product, so its dependence on pure maps and on intermediate
factorizations is part of the algebra. Noninvertible comparison cells are
morphisms of the hom categories, not relations subsequently quotiented away.

Write tensors in execution order: $M\odot_C N$ means first $M$, then $N$.
Ordinary composition of arrows remains right-to-left. A context followed by
content therefore has normal form
\begin{equation}\label{eq:normalform}
 H(A,B)=(e_A\downarrow d_B),\qquad
 h=(P,Q,f:e_AP\longrightarrow d_BQ).
\end{equation}
For two such arrows, the middle content and context must interact. A crossing
is the precise additional datum which makes these normal forms closed under
composition. Without it the tensor is a category of factorizations, not a
category of composable mixed computations.

\begin{theorem}[Unified construction]\label{thm:main}
Let $C$ be a small ordinary category. Let a contextad presentation $(K,e)$ and
a contentad presentation $(L,d)$ have their usual invertible associativity
and unit constraints, while retaining arbitrary forward grade morphisms.
Then the following constructions and recovery statements hold.
\begin{enumerate}[label=\textup{(\roman*)}]
\item Their comma modules $M(A,B)=(e_A\downarrow B)$ and
$N(A,B)=(A\downarrow d_B)$ are pseudomonads on $C$ in $\Biprof$.
\item There is an explicit pseudonatural equivalence
$M\odot_C N\simeq H$, with $H$ as in \eqref{eq:normalform}, including all
noninvertible grade cells and all endpoint transports.
\item A pseudodistributive law $\tau:N\odot_C M\Rightarrow M\odot_C N$
constructs a composite pseudomonad on $M\odot_C N$, hence on $H$ by transport.
Its realization is a bicategory $B_H$ with designated pure pseudofunctor
$J:C\to B_H$, and its relative quintet double category retains the full
endpoint-square data.
\item The original one-sided constructions are recovered by specializing
the other effect to the identity. In a fixed mixed object, the unit-induced
one-sided comparisons are not automatically fully faithful. They are locally
fully faithful in the probability/list instance of Section~\ref{sec:prob}.
\item Marking the generic expansion and contraction arrows in the native
one-sided realizations, their factorizations and chosen lifts recovers the
supplied context and content presentations. The structured input retains
these one-sided marked presentations and their mixed comparison maps. Merely
marking their images in a nonfaithful mixed output cannot recover lost data;
no such claim, or intrinsic recognition from unmarked $(B_H,J)$, is made.
\end{enumerate}
\end{theorem}

The subsequent sections define every part of this statement and prove the
newly needed comparison and concrete instance. Established pseudomonad
composition is used at exactly the point indicated in (iii). The construction
contains ordinary monadic and comonadic computations, forward-parameterized
arrows, display spans and their content duals as specializations of its comma
modules. Mixed opwreaths enter as further algebra objects, not necessarily as
products of two independently existing effects. Factorization recognition is
a property of marked substructures inside it. No claim is made that arbitrary
higher noncartesian structures have a canonical presentation of this form.

Fix universes $\mathcal U\in\mathcal V\in\mathcal W$. Initially $C$ and the
fibres are $\mathcal U$-small. Presheaves valued in $\Cat_{\mathcal U}$ form
a $\mathcal V$-small bicategory, and the corresponding biprofunctor ambient is
formed one level higher when needed. In the set-based example use
$\Set_{\mathcal U}$ as a $\mathcal V$-small base, with the associated increase
of the other levels. There is no category of all sets asserted to be small.

\section{Relative pseudopromonads and their realization}\label{sec:ambient}

Lawler's category-valued profunctor tricategory $\Biprof$ has small
bicategories as objects, pseudofunctors as its profunctor 1-cells,
pseudonatural transformations as 2-cells and modifications as 3-cells
\cite[Section 4.1.1]{Lawler}. We first restrict the objects to ordinary
categories. With our variance convention a 1-cell $C\rightsquigarrow D$ is
\[
 X:C^{\mathrm{op}}\times D\longrightarrow\Cat,
 \qquad
 (X\odot_D Y)(a,c)=\coend{b\in D}X(a,b)\times Y(b,c).
\]
Our 1-cell direction is the reverse of Lawler's: his $D\rightsquigarrow C$
is the displayed $C^{\mathrm{op}}\times D$ pseudofunctor. Our $X\odot Y$
is his $X\circ Y$. We use his concrete hom realization, not an unqualified
transport of a Kleisli universal property through this reversal. The unit is $I_C(a,b)=\disc C(a,b)$.

The subscript on the coend is essential. Its objects have representatives
$[b;x,y]$. Its morphisms include product-category arrows and invertible
balance arrows between
\[
 [b';X(1,u)x,y']\quad\hbox{and}\quad[b;x,Y(u,1)y']
 \qquad (u:b\to b').
\]
They satisfy naturality, unit and composition relations, with the action
constraints inserted when actions are weak. These balance arrows are
invertible; the original arrows in $X(a,b)$ or $Y(b,c)$ need not be.
A functor out of this category is exactly a pseudobalanced cocone. Finite
paths in its small generating graph and the stated relations produce a small
category. Iterated tensors represent cocones balanced in all middle variables.
Their rebracketing equivalences and higher coherence are induced from that
common flattened presentation. The unit equivalences evaluate pure actions.
This is the coend realization of the established tricategory, not a proposal
to declare a family of categories coherent by fiat.

\begin{definition}\label{def:relative}
A relative category-valued pseudopromonad on $C$ is a pseudomonad in
$\Biprof$ on the object $C$: a module $H$, transformations
\[
 m:H\odot_C H\longrightarrow H,\qquad j:I_C\longrightarrow H,
\]
and invertible associativity and left/right unit modifications satisfying
the pseudomonad coherence axioms. Equivalently it is a pseudomonoid in the
endomorphism monoidal bicategory. Invertibility of these structural
modifications imposes no invertibility condition on morphisms inside
$H(a,b)$.
\end{definition}

\begin{proposition}[Concrete realization]\label{prop:realization}
A relative pseudopromonad determines a bicategory $B_H$ and an
identity-on-objects pseudofunctor $J:C\to B_H$, with
\[
 B_H(a,b)=H(a,b),\quad g\circ f=m([f,g]),\quad
 1_a=j(\id_a),\quad J(u)=j(u).
\]
Conversely, such a pair $(B,J)$ determines a relative pseudopromonad. These
constructions recover the hom categories unchanged and recover composition
and pure action up to their specified coherent constraints.
\end{proposition}
\begin{proof}
Composition on 2-cells is $m$ applied to the pair of morphisms. The
associativity and unit modifications evaluated on $[f,g,h]$ and pure identity
representatives give the bicategory associator and unitors. Their equations
are precisely the required pentagon and triangle. Evaluating compatibility
with $j$ on pure maps supplies the compositor and unitor of $J$.

In the converse direction put $H(a,b)=B(a,b)$ and
\[
 H(u,v)(f)=J(v)\circ f\circ J(u)
 \quad (u:a'\to a,\ v:b\to b'),
\]
with a fixed bracketing. Action coherence is assembled from the associator of
$B$ and the constraints of $J$. Send $[f,g]$ to $g\circ f$. For a middle
pure map the required balance is the associator between
$g\circ(J(u)\circ f)$ and $(g\circ J(u))\circ f$. Naturality follows from
whiskering functoriality, and balance composition from the associator and
pseudofunctor axioms. The monoid equations are the same bicategory equations.
The unit sends $u$ to $J(u)$.

Realizing this last module returns the given hom categories and composition
with coherent bracketings. Starting instead from $H$, its original pure action
is equivalent to composition with $j$ by the unit modifications and
pseudonaturality of $m$. The hom comparison functors are identities. This is
the concrete hom-category construction underlying
\cite[Theorem 4.1.3 and Definition 4.1.8]{Lawler} used here.
\end{proof}

One can display the same realization as a category-valued matrix
pseudomonad. On $\Ob C$ use
\[
 (X\star Y)(a,c)=\coprod_bX(a,b)\times Y(b,c),\qquad
 \Delta(a,b)=\begin{cases}1&a=b,\\0&a\ne b.\end{cases}
\]
The maps $H\star H\to H\odot_C H\xrightarrow{m}H$ and
$\Delta\to I_C\xrightarrow{j}H$ are the realized composition and identities.
The pure functor remains marked; forgetting it would forget what the tensor
was relative to.

For clarity, the comparison notion needed here is explicit. Over a functor
$F:C\to D$, a strong operational comparison is a pseudofunctor
$T:B\to B'$ with $T(a)=F(a)$ and an invertible icon
\[
 \theta:TJ\Longrightarrow J'F.
\]
Its components $\theta_u:T(Ju)\Rightarrow J'(Fu)$ are 2-cells with identity
object components, compatible with pure composition and units. For fixed
$F$, a comparison cell can be an icon $\rho:T\Rightarrow T'$ satisfying
$\theta'_u\rho_{Ju}=\theta_u$. These icons need not be invertible. This is a
specified usable class of comparisons; it is not asserted to exhaust all
lax, oplax, pseudo or double transformations in the cited theories.

The same data can be stated in modules. Pull back $H'$ to
$F^*H'(a,b)=H'(Fa,Fb)$. Pure maps act through $F$, and composition is induced
from $m'$ by the map of factorization categories which sends the middle
object $b$ to $Fb$. Together with $I_C\to F^*I_D$, this gives a relative
pseudopromonad. A pure-compatible multiplicative transformation
$H\to F^*H'$, with invertible structural comparison modifications, realizes
the displayed $T,\theta$. Thus varying the base is part of the construction,
not something ruled out by an identity-on-objects convention within a single
instance.

\section{All endpoint squares are retained}\label{sec:squares}

\begin{definition}
For $(B,J)$ over an ordinary category $C$, its relative quintet double
category $\Qrel_C(B,J)$ has objects and vertical arrows those of $C$,
horizontal arrows those of $B$, and squares over $u:a\to a'$ and $v:b\to b'$
given by arbitrary 2-cells
\begin{equation}\label{eq:square}
 \alpha:J(v)\circ h\Longrightarrow h'\circ J(u).
\end{equation}
\end{definition}

Globular squares are canonically identified with $B$-cells using the unit
constraints of $J$ and $B$. They are not literally the same displayed
2-cells if $J$ is nonnormal. One may normalize $J$ or use those canonical
identifications; no normality is silently assumed.

For adjacent squares $\alpha:h\to h'$ over $(u,v)$ and
$\beta:k\to k'$ over $(v,w)$, horizontal composition is
\[
 J(w)kh\xRightarrow{\beta*h}k'J(v)h
       \xRightarrow{k'*\alpha}k'h'J(u),
\]
with the structural reassociations inserted. For a square $\gamma:h'\to h''$
over $(u',v')$, vertical composition is
\[
 J(v'v)h\cong J(v')J(v)h
 \xRightarrow{J(v')*\alpha}J(v')h'J(u)
 \xRightarrow{\gamma*J(u)}h''J(u')J(u)
 \cong h''J(u'u).
\]
Interchange is whiskering functoriality; vertical associativity uses the
compositor axiom for $J$, and horizontal associativity that of $B$. The
binding squares of the companion $J(u)$ are identities of $J(u)$ with the
unit constraints inserted. They satisfy the two companion equations.

\begin{theorem}[Companion bending recovery]\label{thm:bend}
Let $\mathbb D$ be a pseudo double category with vertical category $C$ and a
chosen companion $\widehat u$ for every vertical map. The horizontal
bicategory $B=H(\mathbb D)$ has a canonical pure pseudofunctor $J(u)=\widehat u$.
There is a comparison
$\mathbb D\to\Qrel_C(B,J)$ which is identity on objects and both kinds of
arrows and is a bijection on squares with each fixed boundary. It preserves
pasting up to the canonical structural constraints. No conjoint is required.
\end{theorem}
\begin{proof}
Write $\eta_u$ for the binding square with top $1_a$, bottom $\widehat u$,
left $\id_a$, right $u$, and $\epsilon_u$ for that with top $\widehat u$,
bottom $1_{a'}$, left $u$, right $\id_{a'}$. For a square $\alpha$ with
boundary $(h,h',u,v)$, horizontally paste $\eta_u$, $\alpha$, $\epsilon_v$.
The result, after unit constraints, is the globular cell
$J(v)h\Rightarrow h'J(u)$.

Conversely, begin with that globular cell. Vertically paste the square
$h\Rightarrow J(v)h$ obtained by whiskering $\eta_v$ with $h$, then the
cell, then $h'J(u)\Rightarrow h'$ obtained from $\epsilon_u$ and $h'$.
The resulting boundary is $(u,v)$. In either composite of these operations,
the extra adjacent binding squares cancel by the companion equations. The
operations are therefore mutually inverse without an invertibility
assumption on $\alpha$. The same cancellation at internal boundaries proves
pasting preservation. Constraints between chosen companions of composites
are canonical, and their uniqueness proves the remaining coherence.
\end{proof}

A comparison $(F,T,\theta)$ sends \eqref{eq:square} to the pasting
\[
 J'(Fv)T(h)\cong T(Jv)T(h)\cong T(Jv\circ h)
 \xRightarrow{T\alpha}T(h'\circ Ju)
 \cong T(h')T(Ju)\cong T(h')J'(Fu).
\]
Thus it carries noninvertible endpoint squares as well as globular cells.
Local faithfulness and local fullness are properties of the actual hom
functors $T_{a,b}$, not automatic consequences of having such a comparison.
A local equivalence with a coherently pure-compatible inverse gives the
corresponding double equivalence. These concrete facts do not assert a
blanket triequivalence between every possible notion of presentation map,
double functor and pseudomonad map.

\section{Context and content as presentations of the algebra}\label{sec:presentations}

Let $K:C^{\mathrm{op}}\to\Cat$ be the context indexing pseudofunctor with
extension $e:\int K\to C$. Write $e_AP$ for the extension of grade $P\in K_A$.
For $u:A'\to A$, the chosen cartesian lift has extension
\[
 e_u(P):e_{A'}(u^*P)\longrightarrow e_AP.
\]
A contextad has unit grade $I_A$, discharge $\varepsilon_A:e_AI_A\to A$,
and accumulation $P\otimes R\in K_A$ for $R\in K_{e_AP}$, with comparison
\[
 \delta_{P,R}:e_A(P\otimes R)\longrightarrow e_{e_AP}R.
\]
These maps need not be isomorphisms. The contextad coherence supplies
invertible horizontal associativity and unit constraints.

Dually let $L:C\to\Cat$, with extension $d:\int L\to C$ and cocartesian
lift maps $d_v(Q):d_BQ\to d_{B'}(v_!Q)$. Its contentad unit is $U_B$ with
$\eta_B:B\to d_BU_B$. For $Q\in L_B$ and $S\in L_{d_BQ}$ write its ordered
accumulation as $Q\diamond S$ and its comparison as
\[
 \mu_{Q,S}:d_{d_BQ}S\longrightarrow d_B(Q\diamond S).
\]
The choice $U$ avoids confusing a content unit grade with the pure functor
$J$. Both grade categories retain all their given morphisms.

Define
\begin{align*}
 M(A,B)&=(e_A\downarrow B),&
 M(u,v)(P,f)&=(u^*P,vf e_u(P)),\\
 N(A,B)&=(A\downarrow d_B),&
 N(u,v)(Q,g)&=(v_!Q,d_v(Q)gu).
\end{align*}
A morphism of $M$ is $\alpha:P\to P'$ with $f=f'e_A\alpha$.
A morphism of $N$ is $\beta:Q\to Q'$ with $d_B\beta\,g=g'$.
The endpoint actions on cells are $u^*\alpha$ and $v_!\beta$.
Naturality of lift maps verifies these equations after transport. The action
compositors are the cleavage isomorphisms, whose pentagon and unit equations
are those of the indexed pseudofunctors. Neither action is presumed strict.

The context multiplication sends $(P,f):A\Rightarrow B$ and
$(R,g):B\Rightarrow D$ to grade $P\otimes f^*R$ and body
\begin{equation}\label{eq:ctxcomp}
 e_A(P\otimes f^*R)\xrightarrow{\delta}e_{e_AP}(f^*R)
 \xrightarrow{e_f(R)}e_BR\xrightarrow{g}D.
\end{equation}
Its unit is $(I_A,\varepsilon_A)$; pure $u$ is $(I_A,u\varepsilon_A)$.
The content multiplication sends $(Q,f):A\Rightarrow B$ and
$(S,g):B\Rightarrow D$ to grade $S\diamond g_!Q$ and body
\begin{equation}\label{eq:cntcomp}
 A\xrightarrow{f}d_BQ\xrightarrow{d_g(Q)}d_{d_DS}(g_!Q)
 \xrightarrow{\mu}d_D(S\diamond g_!Q).
\end{equation}
Its unit is $(U_A,\eta_A)$; pure $u$ is $(U_B,\eta_Bu)$.
Composition with a pure middle map agrees with endpoint transport through the
unit constraints. Hence these multiplication functors are pseudobalanced
and descend to the relative tensors. Their associativity and unit
modifications are the native ones. This proves part (i) of
Theorem~\ref{thm:main}, including its global module assertion.

Native context squares over $(u:A\to A',v:B\to B')$ are
\[
 \alpha:P\to u^*P',\qquad vf=f'e_u(P')e_A\alpha.
\]
Native content squares are
\[
 \beta:v_!Q\to Q',\qquad d_{B'}\beta\,d_v(Q)f=f'u.
\]
The pure companions for contextads are supplied in
\cite[Theorem 4.24]{CapucciMyers}. For contentads one checks directly: the
first binding square for $(U_B,\eta_Bu)$ uses the unit transport isomorphism
$u_!U_A\cong U_B$ and naturality of $\eta$; the second uses $\id_{U_B}$.
Their companion equations are the unit and transport equations. Equivalently,
composing a content arrow with a pure output map has grade $v_!Q$, whereas
precomposing its target with a pure input map leaves grade $Q'$, so the bent
cell is exactly the displayed $\beta$. Theorem~\ref{thm:bend} therefore
recovers every native square on both sides with its original cell direction.

\begin{proposition}[Marked presentation recovery]\label{prop:marked}
Mark the context generic arrows
$\kappa_P=(P,\id_{e_AP}):A\Rightarrow e_AP$ with their native squares and
factorizations. Their category, allowing both endpoints to vary, recovers
$\int K$, its source projection and extension. Marking cartesian lifts
recovers its chosen cleavage; marking unit and accumulation factorizations
recovers the remaining supplied contextad structure. The dual statement
holds for the content arrows
$\gamma_Q=(Q,\id_{d_BQ}):d_BQ\Rightarrow B$.
\end{proposition}
\begin{proof}
A square $\kappa_P\to\kappa_{P'}$ over $(u,v)$ is
$\alpha:P\to u^*P'$ and its body equation forces
$v=e_u(P')e_A\alpha$. Its data are therefore exactly a morphism in $\int K$;
the endpoint projections are the projection to $C$ and $e$. Chosen
cartesian squares specify the original cleavage. Every context arrow
$(P,f)$ factors as $\kappa_P$ then $J(f)$ through the unit comparison.
The composite $\kappa_R\circ\kappa_P$ is $(P\otimes R,\delta_{P,R})$,
not generally $\kappa_{P\otimes R}$. Its marked factorization through that
expansion followed by $J(\delta_{P,R})$ retains the comparison $\delta$.
Marking the unit factorization similarly retains $\varepsilon$.

For content, a square $\gamma_Q\to\gamma_{Q'}$ over $(r,v)$ is
$\beta:v_!Q\to Q'$ with $r=d_{B'}\beta\,d_v(Q)$, exactly a morphism in
$\int L$. Its projections recover $d$ and the opfibration. General arrows
factor as $J(f)$ then $\gamma_Q$, and marked generic composites retain
$\mu$ and $\eta$. No unique choice of these markings is inferred from an
unmarked operational object.
\end{proof}

\section{The weak relative tensor is the mixed comma category}\label{sec:bridge}

The main technical bridge must be proved for a pseudocoend: an ordinary
Cat-coend calculation alone would not justify weak endpoint actions.

\begin{theorem}[Pseudocoend normalization]\label{thm:coend}
For functors $e:K\to C$ and $d:L\to C$ there is an equivalence
\[
 \Phi:\coend{Z\in C}(e\downarrow Z)\times(Z\downarrow d)
       \simeq(e\downarrow d).
\]
For the indexed extensions above these equivalences are pseudonatural in
both endpoints, giving $M\odot_C N\simeq H$ as modules, where
\[
 H(u,v)(P,Q,f)=(u^*P,v_!Q,d_v(Q)f e_u(P)).
\]
They retain arbitrary morphisms of $K$ and $L$.
\end{theorem}
\begin{proof}
Present the pseudocoend $T$ by objects
$t=[Z;P,x:eP\to Z;Q,y:Z\to dQ]$, product arrows $(\alpha,\beta)$ with
$x=x'e\alpha$ and $y'=d\beta\,y$, and invertible balance generators
\[
 b_u:[Z;P,x;Q,y'u]\longrightarrow[Z';P,ux;Q,y']
 \qquad(u:Z\to Z').
\]
Impose product functoriality, naturality of $b$, $b_{\id}=\id$ and
$b_{vu}=b_vb_u$ with the evident arguments. This is the universal
pseudobalanced-cocone presentation. No grade arrow is declared invertible.

Define $\Phi(t)=(P,Q,yx)$, $\Phi(\alpha,\beta)=(\alpha,\beta)$, and
$\Phi(b_u)=\id$. All relations hold by composition in $C$. Define an inverse
normalization on objects by
\[
 \Psi(P,Q,f)=[eP;P,\id;Q,f].
\]
For a comma morphism $(\alpha,\beta)$ satisfying
$d\beta\,f=f'e\alpha$, use the composite
\begin{align*}
 [eP;P,\id;Q,f]&\xrightarrow{\beta}
 [eP;P,\id;Q',d\beta f]\\
 &\xrightarrow{b_{e\alpha}}
 [eP';P,e\alpha;Q',f']
 \xrightarrow{\alpha}[eP';P',\id;Q',f'].
\end{align*}
Identities use $b_{\id}$. For two successive comma morphisms, balancing
naturality moves the second output-grade update past the first balance;
input and output updates commute in the product category; and the two
balances compose to $b_{e(\alpha_2\alpha_1)}$. This proves functoriality of
$\Psi$ rather than merely defining it on objects.

Now $\Phi\Psi=\id$ on objects and arrows. At $t$ the generator
\[
 b_x:\Psi\Phi(t)=[eP;P,\id;Q,yx]\longrightarrow t
\]
is invertible. Its naturality for a balance generator is
$b_{ux}=b_ub_x$. For a product arrow it follows from
$x=x'e\alpha$, $b_x=b_{x'}b_{e\alpha}$, and naturality in both grade
variables. These generators cover every arrow of $T$, so $b_x$ defines a
natural isomorphism $\Psi\Phi\Rightarrow\id$. This proves the equivalence
without either identifying balance arrows literally or discarding grade
arrows.

Apply this at $e_A,d_B$. The endpoint functor on $H$ is the one in the
statement, and on morphisms is $(u^*\alpha,v_!\beta)$. Naturality of lifts
proves its body equation. Both routes through $\Phi$ have composite body
$d_v(Q)yx e_u(P)$, with exactly the cleavage constraints for weak choices.
For $\Psi$, source transport changes the normal splitting object from
$e_AP$ to $e_{A'}u^*P$. The comparison is the specific generator
$b_{e_u(P)}$, not an arbitrary choice of representative equivalence.
Output transport keeps the splitting object and uses the content cleavage
constraint. Iterated source transports satisfy their comparison equation
because lift composition and $b_{vu}=b_vb_u$ do; mixed source/output
transports commute by balance naturality. Thus the equivalence is global
and pseudonatural.
\end{proof}

If balances are imposed as equalities instead, the corresponding ordinary
strict Cat-coend is isomorphic to the comma category by the same
normalization. This does not identify arbitrary strict and pseudo coends;
the equivalence just proved uses the special factorization through $eP$.

In coordinates a mixed square from $(P,Q,f):A\Rightarrow B$ to
$(P',Q',f'):A'\Rightarrow B'$ over $(u,v)$ is
\begin{equation}\label{eq:mixedsquare}
 \alpha:P\to u^*P',\quad \beta:v_!Q\to Q',\qquad
 d_{B'}\beta\,d_v(Q)f=f'e_u(P')e_A\alpha.
\end{equation}
It is a morphism $H(1,v)(h)\to H(u,1)(h')$ in $H(A,B')$.
Once multiplication is installed, the unit constraints identify these
actions with pure whiskering, so \eqref{eq:mixedsquare} is precisely the
relative-quintet square. Its vertical composition uses the actual endpoint
functors. Its horizontal composition comes from multiplication, without an
isomorphism-class quotient.

\section{One mixed multiplication from a coherent crossing}\label{sec:mixed}

\begin{definition}\label{def:law}
An interaction between $M$ and $N$ is a pseudodistributive law
$\tau:N\odot_C M\Rightarrow M\odot_C N$. Its crossing need not be
invertible. It has invertible modifications, suppressing brackets,
\begin{align*}
 \tau(j_NM)&\cong Mj_N,& \tau(Nj_M)&\cong j_MN,\\
 \tau(\mu_NM)&\cong(M\mu_N)(\tau N)(N\tau),&
 \tau(N\mu_M)&\cong(\mu_MN)(M\tau)(\tau M),
\end{align*}
subject to the complete pseudodistributive-law coherence conditions, not just
these four comparison boundaries. Precisely, we use the Gray-category
definition of \cite[Section 4]{Marmolejo}, with its nine coherence equations,
transported through a Gray presentation of $\Biprof$. The equivalent modern
formulation is \cite[Definition 4.1]{GambinoLobbia}, including its axioms
\textup{(C1)--(C8)}. The associators and interchangers of the ambient are
included in these equations.
\end{definition}

The reference to the full standard law is an input hypothesis: no generic
algorithm below produces it from arbitrary $K,L$. In the main example all
its coherence equations are verified by an address-preserving construction,
so they are not an untested assumption there.

\begin{theorem}[Mixed algebra]\label{thm:mixed}
Given such an interaction, $M\odot_C N$ carries the unit and multiplication
\begin{align}
 j_{MN}&:I_C\simeq I_C\odot_C I_C
       \xrightarrow{j_M\odot j_N}M\odot_C N,\label{eq:unitMN}\\
 \mu_{MN}&:MNMN\xrightarrow{M\tau N}MMNN
            \xrightarrow{\mu_M\mu_N}MN.\label{eq:multMN}
\end{align}
Transport along Theorem~\ref{thm:coend} makes $H(A,B)=(e_A\downarrow d_B)$
a relative pseudopromonad. Its realization and relative quintets give a
single mixed structure with the endpoint squares \eqref{eq:mixedsquare}.
\end{theorem}
\begin{proof}
The composite-pseudomonad theorem is
\cite[Theorem 5.1, pp.\ 101--102]{Marmolejo}. Its construction is explicitly
$DUDU\to DDUU\to DU$, with unit and associator pastings supplied there; its
hypothesis allows an arbitrary crossing 2-cell with invertible structural
3-cells. Every tricategory has a Gray presentation, so the theorem applies
there and transports back to $\Biprof$. The tensor $MN$ is identified with
the concrete pseudocoend tensor by coherent equivalence; it is not replaced
by an unspecified new carrier. Theorem~\ref{thm:coend} then transports the
unit, multiplication and coherence along the explicitly given equivalence
$MN\simeq H$. Evaluating these operations realizes the bicategory, and
Section~\ref{sec:squares} supplies its squares.

Concretely the associativity comparison evaluates the two normalization
routes for $MNMNMN$. The two multiplication modifications of $\tau$ and
the input associators compare the routes; their full law coherence gives
the pentagon, while the unit modifications give the triangle. Thus
associativity on $H$ is inherited from the law, not separately imposed after
an arbitrary crossing has been chosen. The lifting theorem
\cite[Theorem 4.5]{GambinoLobbia} alone is not the composition theorem being
used here.
\end{proof}

The formula can be evaluated directly on grades. At the generic pair
\[
 (Q,\id):d_BQ\Rightarrow B,\qquad(R,\id):B\Rightarrow e_BR,
\]
normalize the value of $\tau$ to obtain
\[
 \widehat R\in K_{d_BQ},\qquad \widehat Q\in L_{e_BR},\qquad
 \chi:e_{d_BQ}\widehat R\longrightarrow d_{e_BR}\widehat Q.
\]
For $h=(P,Q,f):A\Rightarrow B$ and $k=(R,S,g):B\Rightarrow D$, the composite
has grades
\begin{equation}\label{eq:mixedgrades}
 P\otimes f^*\widehat R,\qquad S\diamond g_!\widehat Q
\end{equation}
and body
\begin{align}\label{eq:mixedbody}
 e_A(P\otimes f^*\widehat R)
 &\xrightarrow{\delta}e_{e_AP}(f^*\widehat R)
 \xrightarrow{e_f(\widehat R)}e_{d_BQ}\widehat R
 \xrightarrow{\chi}d_{e_BR}\widehat Q\notag\\
 &\xrightarrow{d_g(\widehat Q)}d_{d_DS}(g_!\widehat Q)
 \xrightarrow{\mu}d_D(S\diamond g_!\widehat Q).
\end{align}
Every displayed arrow is typed. Functoriality on cells comes from $\tau$,
cartesian/cocartesian transport and the input multiplications. The pure arrow
$u:A\to B$ is represented by
\[
 (I_A,U_B,\eta_Bu\varepsilon_A).
\]
With the other side specialized to $I_C$, pseudoco-Yoneda gives $M\odot I_C
\simeq M$ or $I_C\odot N\simeq N$. In a fixed mixed object the unit-induced
maps are
\begin{equation}\label{eq:onesided}
 (P,f)\longmapsto(P,U_B,\eta_Bf),\qquad
 (Q,g)\longmapsto(I_A,Q,g\varepsilon_A).
\end{equation}
The law's unit equations make these pure-compatible strong comparisons.
Their hom functors may identify bodies or introduce additional cells. For
example a nonmonic $\eta_B$ can identify distinct context bodies; with the
terminal monad on $\Set$ every such body is identified. Full recovery in a
fixed mixed object requires a separate hom-category argument.

A candidate interaction must first descend through pure-middle balance.
For example the previously tested finite image/pushout recipe gives, for two
balanced representatives, enabled domains with two elements and one element
respectively. Those domains are not isomorphic over the source. Consequently
that recipe is not a functor from the pseudo-balanced tensor into a target
retaining those domains, even before its associativity is considered. The
relative tensor makes this obstruction visible at the correct stage.

\section{Dependent probability and lists in one nonthin object}\label{sec:prob}

Let $\Prob$ have finite nonempty probability spaces with strictly positive
weights and deterministic measure-preserving maps. A morphism $a:P\to Q$
satisfies $q(z)=\sum_{x:a(x)=z}p(x)$. This category has noninvertible maps and
distinct parallel maps; it is not replaced by its groupoid core. Work over
$C=\Set_{\mathcal U}$ at the universe level specified above. Set
\[
 K_A={\Prob}^A,\qquad e_AP=\coprod_{a\in A}P_a.
\]
Reindexing is substitution along functions. The unit grade is the singleton
family. Accumulation of $P$ and $R$ has fibre
$\sum_{x\in P_a}R_{(a,x)}$, with weight $p_a(x)r_{a,x}(y)$; its extension
comparison is the corresponding tuple reassociation. For content take
terminal grade fibres and $d_B(*)=\List B$, with singleton insertion and
concatenation. The resulting mixed hom category is
\[
 H(A,B)=\{(P,f:\textstyle\coprod_aP_a\to\List B)\},
\]
with a morphism $\alpha:(P,f)\to(P',f')$ precisely a family of
measure-preserving maps satisfying $f'(a,\alpha_a x)=f(a,x)$.
Endpoint squares over $(u,v)$ have
\begin{equation}\label{eq:probsquare}
 \alpha_a:P_a\to P'_{u(a)},\qquad
 f'(u(a),\alpha_ax)=\List(v)(f(a,x)).
\end{equation}
There is no output-distribution quotient in this definition.

For a list $\ell=[b_1,\ldots,b_n]$ and grade family $R\in K_B$, define
\begin{equation}\label{eq:probchi}
 \widehat R_\ell=\prod_{i=1}^nR_{b_i},\qquad
 \widehat r_\ell(r_1,\ldots,r_n)=\prod_{i=1}^nr_{b_i}(r_i),\qquad
 \chi(\ell,(r_i))=[(b_i,r_i)]_{i=1}^n.
\end{equation}
The notation for weights is understood as the weight function in the
corresponding fibre. Every \emph{occurrence} has its own coordinate, even
when labels repeat. The empty product is a singleton. This is independent
probability tensor, not a claim that this carrier product is the categorical
product in $\Prob$.

\begin{theorem}[The probability and list law]\label{thm:problaw}
The crossing \eqref{eq:probchi} defines a pseudodistributive law for the
probability-context pseudomonad and list-content pseudomonad. Its mixed
composition is the following explicit operation. If
$h=(P,f):A\Rightarrow B$, $k=(R,g):B\Rightarrow D$, and
$f(a,x)=[b_1,\ldots,b_n]$, then
\begin{align}
 (k\circ h).\mathrm{grade}_a
   &=\sum_{x\in P_a}\prod_{i=1}^nR_{b_i},\label{eq:probgrade}\\
 \mathrm{weight}(x,(r_i))&=p_a(x)\prod_i r_{b_i}(r_i),\notag\\
 (k\circ h)(a,x,(r_i))&=\mathop{\mathrm{concat}}_i g(b_i,r_i).
 \label{eq:proboutput}
\end{align}
All native grade maps and endpoint squares remain present. Both maps in
\eqref{eq:onesided} are locally fully faithful in this instance.
\end{theorem}
\begin{proof}
First, all fibres in \eqref{eq:probgrade} are finite nonempty full-support
probability spaces. Positivity is immediate. Summing the weights gives
$\sum_xp_a(x)\prod_i\sum_{r_i}r_{b_i}(r_i)=1$, including empty lists.
For grade maps $\beta_b:R_b\to R'_b$ the product map preserves probability:
its mass over $(r'_i)$ is
\[
 \sum_{(r_i):\,\beta_{b_i}r_i=r'_i}\prod_i r_{b_i}(r_i)
 =\prod_i\sum_{r_i:\,\beta_{b_i}r_i=r'_i}r_{b_i}(r_i)
 =\prod_i r'_{b_i}(r'_i).
\]
It preserves the list output coordinatewise. Under pure $v:B\to B'$ a
position labelled $b_i$ becomes the position labelled $v(b_i)$; positions
are not merged. Thus the crossing respects pure balance even for
noninjective $v$. These statements define the crossing on all cells and its
pseudobalanced descent, not only its action on isomorphism classes.

For parallel mixed cells $\alpha,\beta$, horizontal composition sends
\[
 (x,(r_i))\longmapsto(\alpha_a x,(\beta_{b_i}r_i)).
\]
Their output equation ensures the target has the same occurrence positions.
For each target point, sum first over the product preimage of the $r'_i$,
then over the $\alpha$-preimage of $x'$. The preceding product identity and
measure preservation of $\alpha$ give exactly the target weight. The same
calculation with \eqref{eq:probsquare} proves compatibility with arbitrary
endpoint squares.

The two multiplication modifications of the law have concrete carriers.
For list multiplication, product over a concatenated list is identified
with nested products over its ordered blocks. For context multiplication,
the comparison is the bijection
\begin{equation}\label{eq:probdistrib}
 \prod_i\left(\sum_{y_i\in R_i}S_{i,y_i}\right)
 \longleftrightarrow
 \sum_{(y_i)\in\prod_iR_i}\prod_iS_{i,y_i}.
\end{equation}
Both sides assign weight $\prod_i r_i(y_i)s_{i,y_i}(z_i)$ to the same
coordinates. Unit comparisons insert or delete singleton factors; empty
list products remain singleton spaces. Every comparison is therefore an
invertible, measure-preserving grade map, natural in noninvertible grade
maps as well.

To check the full law coherence, label a coordinate by its finite sequence
of list positions and context stages. Each structural comparison just
changes bracketing or grouping of that finite addressed family. It leaves
both the value at every address and its factor in the product weight
unchanged. Every composite in each law-coherence diagram is therefore the
same function between the same addressed carriers. A diagram involving a
unit deletes only a forced singleton coordinate; the statement remains
true when a list block is empty. This verifies all coherence diagrams of
Definition~\ref{def:law}, not merely the four object-level boundaries.

For completeness the resulting associator for three computations has the
particularly transparent form
\begin{equation}\label{eq:probassoc}
 ((x,(y_i)),(z_{ij})_{\mathrm{flattened}})
 \longleftrightarrow
 (x,((y_i,(z_{ij})_j))_i).
\end{equation}
The output of the second computation specifies the length of each block,
so variable list lengths do not cause an ambiguous splitting. The map
preserves every addressed sample and the product weight. All pentagon
routes for four computations preserve the same addresses, hence coincide;
triangles delete singleton samples. This also proves the operational
coherence directly.

Finally the context inclusion uses singleton output lists. Its mixed
cell equation is $[f'(a,\alpha x)]=[f(a,x)]$, exactly the native context
cell equation. Thus it is full and faithful on each hom category.
The content inclusion uses singleton sample families. There is exactly one
map between any two such families, and it satisfies the mixed equation
exactly when the two list functions agree. Its hom category is therefore
the discrete list-Kleisli hom category. The unit/associator tuple maps agree
with the native horizontal structures. Endpoint recovery follows either
from these same equations or Theorem~\ref{thm:bend}.
\end{proof}

This instance is genuinely nonthin. For a constant list output, both
permutations of an equally weighted two-element sample space are distinct
parallel cells, and the map from that sample space to a singleton is a
noninvertible cell. All three survive in $H$. Pushing presentations forward
to distributions on lists forgets them and is a further semantics map, not
the mixed construction itself. Finite regression tests can check examples,
but the proofs above hold for every finite family, including repeated
labels, zero-length lists and changing endpoints.

\section{Internal constructors and familiar specializations}\label{sec:constructors}

The comma presentations are flexible because they retain a grade category
rather than demanding a single representing endofunctor. Nevertheless the
familiar fixed-effect forms are literal specializations. A monad $T$ gives
terminal content fibres and $d_B(*)=TB$, hence
\[
 H_T(A,B)=\disc C(A,TB),\qquad g\circ f=\mu_D T(g)f,
 \qquad J(u)=\eta_Bu.
\]
Its endpoint equation is $T(v)f=f'u$. A comonad $G$ gives terminal context
fibres and $e_A(*)=GA$, hence
\[
 H_G(A,B)=\disc C(GA,B),\qquad g\circ f=gG(f)\delta_A,
 \qquad J(u)=u\varepsilon_A,
\]
with endpoint equation $vf=f'G(u)$. These are literal operational and square
recoveries; their discreteness is native to an ordinary base, not a loss of
higher input data. Singleton grade fibres do not make the structural maps
$\eta,\mu,\varepsilon,\delta$ invertible.

For a separate mixed law $\lambda:GT\Rightarrow TG$ the same factory is
\begin{equation}\label{eq:fixedmixed}
 H(A,B)=\disc C(GA,TB),\quad
 g\circ f=\mu_D T(g)\lambda_B G(f)\delta_A,\quad
 J(u)=\eta_Bu\varepsilon_A.
\end{equation}
The endpoint equation is $T(v)f=f'G(u)$. The comparison maps are $f\mapsto
\eta f$ and $g\mapsto g\varepsilon$, with exactly the possible failures of
faithfulness described above.

For a right monoidal action of $V$ on $C$, use $K_A=V$ with constant
reindexing and $e_AP=A\mathbin{\odot}P$. The resulting hom objects are
$(P,f:A\odot P\to B)$, and their cells satisfy
$f=f'(1_A\odot\alpha)$. Composition has parameter $P\otimes Q$ and body
\[
 A\odot(P\otimes Q)\cong(A\odot P)\odot Q
 \xrightarrow{f\odot1}B\odot Q\xrightarrow{g}D.
\]
Endpoint squares are $vf=f'(u\odot\alpha)$. This is the
forward-reparameterization convention for Para. A source using the opposite
cell convention must be opposited accordingly. No copying operation on a
general monoidal parameter has been introduced.

If $\mathcal D$ is a pullback-stable class of maps containing identities and
closed under composition, choose $K_A=\mathcal D/A$, extension the domain,
and reindexing by chosen pullbacks. The same comma hom is the category of
spans $A\xleftarrow{p}X\xrightarrow{f}B$ with $p\in\mathcal D$ and all
commuting apex maps. Composition has apex $X\times_BY$. Endpoint squares are
apex maps $a:X\to X'$ satisfying $p'a=up$ and $f'a=vf$. Thus both apex maps
and changing-endpoint squares survive. Chosen pullback coherence is
canonical, not assumed strict. Pushout-stable codisplays give the dual
content/cospan instance. Having both does not supply a mixed crossing.

Reader and writer transpositions provide controls on what an exact recovery
can mean. For $GA=E\times A$, currying identifies coKleisli arrows with
$A\to B^E$, and both compositions use the same environment twice. The
endpoint equations coincide, giving a pure-square double-category
isomorphism when the requisite exponentials exist. For a monoid $W$ the
writer monad $TA=W\times A$ can alternatively be presented by the discrete
context fibre of functions $w:A\to W$, extension $A$, and composition
\[
 (w,f)\then(v,g)=\bigl(a\mapsto w(a)v(f(a)),\,gf\bigr).
\]
Its grade condition over $u$ is $w=w'u$, exactly the writer endpoint
equation together with the body equation. This is a concrete double
isomorphism, not an asserted equivalence of arbitrary nonthin Para with
reader or writer. The fixed reader/writer law is simply
$(e,(w,x))\mapsto(w,(e,x))$.

\subsection{Grothendieck extension is a base changing constructor}

Assume the split indexed presentation and coherence hypotheses used by the
Grothendieck extension construction. On the 2-category of category-valued
presheaves define
\begin{equation}\label{eq:groth}
 (\mathcal P_KY)(A)=\int_{P\in K_A}Y(e_AP).
\end{equation}
Here the integral is the contravariant Grothendieck construction: a morphism
$(P,y)\to(P',y')$ is
\[
 (\alpha:P\to P',\;\gamma:y\to Y(e_A\alpha)(y')).
\]
Reindex along $u:A\to B$ by
$(P,y)\mapsto(u^*P,Y(e_u(P))y)$. The unit and multiplication are
\begin{equation}\label{eq:grothmonad}
 y\longmapsto(I_A,Y(\varepsilon_A)y),\qquad
 (P,(R,y))\longmapsto(P\otimes R,Y(\delta_{P,R})y).
\end{equation}
These are well typed without invertibility of $\varepsilon$ or $\delta$.
Their action on arrows is induced by the Grothendieck construction and the
native naturality maps. The contextad coherence gives a pseudomonad.
Structurally, the extension is the monoidal pseudofunctor of
\cite[Construction 3.1 and Proposition 3.2]{Kamiya}. This construction is
available before the extra comprehension-equivalence assumption of that
paper's Section 3.2.

\begin{proposition}[Representable realization]\label{prop:groth}
The restriction of $\Kl(\mathcal P_K)$ to representable presheaves recovers
the context operational bicategory and its endpoint squares by a local
equivalence compatible with composition. The entire presheaf Kleisli
bicategory is a further algebra object at its larger bicategorical base,
not an equivalent replacement for that restriction.
\end{proposition}
\begin{proof}
Write $yA=C(-,A)$ as a discrete-valued representable. Pseudonatural Yoneda
and \eqref{eq:groth} give
\[
 \Ps(yA,\mathcal P_KyB)\simeq(\mathcal P_KyB)(A)
                 =(e_A\downarrow B)=M(A,B).
\]
On morphisms, the component $\gamma$ in the discrete representable exists
exactly when $f=f'e_A\alpha$, so this equivalence retains grade maps rather
than only objects. Evaluating Kleisli multiplication using
\eqref{eq:grothmonad} produces exactly \eqref{eq:ctxcomp}; unit evaluation
produces discharge. Yoneda on pure maps supplies the base-changing pure
comparison, hence the endpoint squares are the same ones under bending.
Arbitrary presheaves are additional objects of the full Kleisli bicategory.
Their presence does not imply they are equivalent to representables.

That full construction uses Lawler's full bicategorical-base ambient:
$\Psh(C)$ is a 2-category, not an ordinary pure category. If $S$ is its
pseudomonad, the corresponding module has hom category
$\Psh(C)(X,SY)$ and Kleisli composition. The ordinary-base quintet result
used elsewhere is the representable restriction, or a separately specified
ordinary tight subcategory. This explicitly separates the two base levels.
\end{proof}

For content use the covariant construction on $\Copsh(C)=[C,\Cat]$,
\[
 (\mathcal Q_LY)(B)=\int_{Q\in L_B}Y(d_BQ),\qquad
 (\beta,\gamma):(Q,y)\to(Q',y'),\quad
 \gamma:Y(d_B\beta)y\to y'.
\]
Insertion and multiplication use $Y\eta$ and $Y\mu$. For
$y^A=C(A,-)$ one gets
\[
 \Ps(y^B,\mathcal Q_Ly^A)\simeq(\mathcal Q_Ly^A)(B)
                        =(A\downarrow d_B).
\]
Because $A\mapsto y^A$ is contravariant, content is recovered from the
\emph{1-cell-opposite} of the corepresentable Kleisli restriction. Morphisms
inside the hom categories are not reversed. This is the explicit covariant
dual construction, not an attribution of an unprinted formula to Kamiya.

\subsection{Telescopes evaluate inside the same constructor}

For a split left-fibrant span $V$ in the setting of the extension theorem,
let $V^{\ltimes n}$ be its $n$-fold sequential product. Define
\[
 \Tel(V)=\coprod_{n\geq0}V^{\ltimes n}.
\]
The telescope-monad unit inserts a length-one string, and multiplication
flattens strings of strings. Each fibre is a category of \emph{dependent}
finite chains: the next grade lies over the extension by previous grades,
not in one constant fibre. Monoidality and coproduct preservation give
\begin{equation}\label{eq:telescope}
 \mathcal P_{\Tel(V)}\simeq\coprod_{n\geq0}\mathcal P_V^n.
\end{equation}
The right side is the free-monad construction in the applicable
coproduct-preserving endofunctor setting. Here the free monad's unit is the
length-zero summand, distinct from the telescope-monad unit at $V$ just
described. A contextad multiplication evaluates finite chains by iterated
accumulation, and its coherence makes evaluation compatible with flattening.
Strict telescope algebras recover strict contextads under the specified
normal-cleavage and strict-map assumptions
\cite[Corollary 2.20]{Kamiya}; weak or transported variants require their
coherent analogues. Identifying arbitrary indexed data with a prescribed
display class is stronger and still requires the comprehension equivalence.

\subsection{Opwreaths and factorization are internal structures}

A mixed opwreath need not split into a separate comonad and monad. Street's
input is a monad $T$, an endofunctor $G$, and
\[
 \zeta:GT\Rightarrow TG,\qquad d:G\Rightarrow TGG,\qquad e:G\Rightarrow T
\]
satisfying the opwreath axioms \cite[Definition 5 and Proposition 7]{Street}.
It gives the same-shaped hom sets $C(GA,TB)$ with unit $e_A$ and composite
\[
 GA\xrightarrow{d}TGGA\xrightarrow{TGf}TGTB
 \xrightarrow{T\zeta_B}TTGB\xrightarrow{TTg}TTTD
 \xrightarrow{\mu^{(3)}}TD.
\]
Naturality gives pure maps $J(u)=T(u)e_A=e_BG(u)$, so this is directly an
object of the relative construction. Its theorem supplies associativity;
it need not be the mixed product of a pre-existing coKleisli category, since
$G$ need not be a comonad. Higher two-sided Kleisli versions require their
specified bicategorical input and coherence \cite{Garner}; they are not
claimed to follow from Set-indexed formulas without those data.

Factorization recognition instead marks classes or subbicategories inside
$(B,J)$ and asks whether their composition has the stated universal
factorization property. In strict discrete examples unique factorization
gives a normal form and hence a distributive crossing; with cells the
appropriate datum is an equivalence of factorization categories with
coherence. Rosebrugh--Wood and Lack develop such reconstruction under their
factorization hypotheses \cite{RosebrughWood,LackProps}. The double-category
recognition results of \v{S}t\v{e}p\'an concern particular filling,
bicartesian and uniqueness conditions
\cite[Theorems 3.3.7 and 3.3.33]{Stepan}. They identify structures already
expressible here, not an extra universal effect that automatically commutes
with the others. For a precise envelope, let $A$ be ordinary with wide subcategories
$\mathcal E,\mathcal M$. Regard $A$ as locally discrete and form
$\Qrel_{\mathcal E}(A,J_{\mathcal E})$ with tight arrows $\mathcal E$.
Mark and restrict its horizontal arrows to $\mathcal M$. Its squares are
exactly commuting squares $vf=gu$, which recovers the ordinary
factorization double category literally. The $\mathcal E$-companions exist
in the envelope but need not belong to the restricted $\mathcal M$ part.
Thus a factorization double need not itself have companions. Under the
cited recognition hypotheses one first reconstructs $A,\mathcal E,\mathcal M$
and then applies this envelope; arbitrary nonflat or companionless double
categories are not thereby automatically covered. A factorization theorem
can supply a crossing; a quotient that destroys factorization witnesses can
also destroy its hypotheses.

\section{Comparison maps and controlled forgetting}\label{sec:comparisons}

The construction operates on compatible packages, not only on objects.
Over a fixed base, let $F:M\to M'$ and $G:N\to N'$ be strong relative
pseudomonad maps. Suppose an invertible modification identifies
\[
 (F\odot G)\tau\quad\hbox{with}\quad\tau'(G\odot F),
\]
compatibly with the four law modifications and all their coherence.
Then $F\odot G:MN\to M'N'$ preserves the mixed units and multiplication
through the specified constraints. Indeed, substituting the displayed
modification into $M\tau N$ and then the multiplicative comparisons of
$F,G$ produces the composite multiplication comparison; its associativity
and unit equations are the assumed law-map equations. Transporting through
the comma equivalences gives the corresponding map on mixed presentations.
This is one precise class of factory maps, not a claim that all competing
transformation conventions coincide.

There is also a useful, explicitly separate lax comparison class. Specify a
lax double functor directly between the relative quintets: a functor on the
vertical bases, assignments on horizontal arrows and \emph{all} squares,
and natural comparisons
\[
 T(k)\circ T(h)\Longrightarrow T(k\circ h),\qquad
 1_{T(a)}\Longrightarrow T(1_a),
\]
with the double-functor unit, associativity and square-pasting conditions.
These comparisons need not be invertible. Merely replacing a pseudofunctor
by a lax functor in Section~\ref{sec:squares} is insufficient: the square
transport there uses the inverse compositor at $(h',J(u))$. One must
supply the square action directly or prove the required mixed-pure
compositors invertible. This enlarged comparison class is not identified
with ordinary 3-cells of $\Biprof$.

For example, over $\mathbf{FinSet}$ let $\mathrm{HistProb}(A,B)$ be the
probability-context hom category with deterministic output $h_a:P_a\to B$.
Composition keeps both samples. Let $\mathrm{FinStoch}(A,B)$ be the
discrete category of finite stochastic kernels. The map
\begin{equation}\label{eq:historyE}
 (Eh)_a(b)=\sum_{x:h_a(x)=b}p_a(x)
\end{equation}
is a strong double comparison. Finite summation proves preservation of
convolution; every history square over $(u,v)$ gives the kernel square
$v_*Eh=(Eh')u$.

A converse support-sampler assignment $S$ sends a kernel $k$ to sample space
$\operatorname{supp}(k_a)$ with weight $k_a$ and output its inclusion into
$B$. For a kernel square $v_*k=k'u$, its history-square map is
\[
 \operatorname{supp}(k_a)\longrightarrow
 \operatorname{supp}(k'_{u(a)}),\qquad b\longmapsto v(b).
\]
It is well-defined because positive mass stays positive in pushforward;
the kernel equation says exactly that this map preserves probability.
It preserves vertical square composition. The horizontal comparison
$S\ell\circ Sk\Rightarrow S(\ell\circ k)$ sends a retained pair $(b,c)$
to $c$. Convolution proves measure preservation; both associativity paths
send $(b,c,d)$ to $d$; unit paths delete singleton coordinates. Naturality
against an arbitrary endpoint square sends $(b,c)$ to $(v(b),w(c))$, and
both routes then give $w(c)$. Thus $S$ is a direct lax double comparison,
not merely an object-level sampling recipe.

\begin{proposition}[A local reflection, not an equivalence]
At every pair of endpoints, $E$ is left adjoint to $S$, with $ES=\id$.
The lax comparison of $S$ need not be invertible.
\end{proposition}
\begin{proof}
A history morphism $h\to Sk$ must send every sample $x$ to its output
$h_a(x)$. It is therefore unique if $Eh=k$, and nonexistent otherwise,
which is exactly the hom set from $Eh$ to $k$ in the discrete kernel
category. The unit is the output map on samples and the counit is identity.
A fair two-point kernel followed by a constant kernel has two retained
samples in $S\ell\circ Sk$, whereas $S(\ell\circ k)$ has one. The
comparison collapses two points and is not invertible.
\end{proof}

This illustrates why retaining noninvertible cells and permitting controlled
forgetting are distinct. The common construction can express both, with the
correct comparison strength, without presenting stochastic semantics as
an equivalence of history presentations. For base-changing strong double
comparisons, vertical transformations may likewise be retained directly:
a natural transformation $t:F\Rightarrow G$ of base functors together with
squares $J'(t_b)T(h)\Rightarrow T'(h)J'(t_a)$, natural against every source
square and compatible with unit and composition comparisons. Their cells
need not be invertible; no blanket identification with all pseudomonad
transformation notions is required.

\section{The exact obstruction to universal conservative assembly}\label{sec:obstruction}

Three assertions must remain distinct. A common ambient holds two algebra
objects. A common computation object has maps from both sides agreeing on
the pure interface. A conservative mixed normal form additionally has a
specified short factorization category, lawful multiplication and recovery
properties for those maps. The first assertion implies neither the second
with faithful maps nor the third. Even faithful pure interfaces do not
remove the obstruction.

\begin{theorem}[Bicyclic obstruction]\label{thm:bicyclic}
There are faithful functors from a common ordinary category $C$ into
categories $M,N$ such that no category receiving both, with their pure maps
identified and their equations preserved, can receive both faithfully.
For locally discrete bicategories the analogous common target cannot make
both inclusions locally fully faithful, even if pure agreement is only
specified up to coherent isomorphism.
\end{theorem}
\begin{proof}
Use one-object categories, writing multiplication as ordinary function
composition. Let
\[
 C=\langle a\rangle,\qquad
 M=\langle a,b\mid ab=1\rangle,\qquad
 N=\langle a,c\mid ca=1\rangle.
\]
The pure functors are faithful. For $M$, represent $a(n)=\max(n-1,0)$ and
$b(n)=n+1$ on natural numbers. Distinct powers of $a$ are distinct functions,
$ab=1$, and $ba\ne1$ since $ba(0)=1$. For $N$, use $a(n)=n+1$ and
$c(n)=\max(n-1,0)$; again powers of $a$ are distinct, $ca=1$ and $ac\ne1$.
In any common ordinary target,
\[
 c=c(ab)=(ca)b=b,
 \qquad ba=ca=1.
\]
Thus $M$ cannot map faithfully there; symmetrically $ac$ is collapsed in
$N$. In fact the ordinary pushout is the infinite cyclic group, though this
identification is unnecessary for the obstruction.

For bicategorical targets, transport along the coherent comparison of the
two pure images. The image of $a$ has both a right inverse $b$ and a left
inverse $c$ up to invertible 2-cells. The same calculation, with associators
inserted, gives $b\cong c$ and $ba\cong1$. But the discrete source hom
category of $M$ has no morphism $ba\to1$, so its image cannot be locally
fully faithful. This is a failure of full recovery, not a claim of failure
of local faithfulness on 2-cells, which would be vacuous in a discrete hom
category.
\end{proof}

The example is an obstruction to a generic common-base assembly theorem;
it does not assert that these particular monoids arise from the specific
dependent presentations used above. It also does not prohibit a
nonconservative mixed output. That is why Theorem~\ref{thm:mixed} separates
existence of law data from full faithfulness of its unit-induced maps.

If one instead wants universal alternating syntax, a safe strict
construction is available. For small strict 2-categories and strict functors
from a common strict 2-category, take the object pushout, generate 1- and
2-cells from both sides with shared generators identified, impose each
side's identity/composition equations, and close under strict 2-category
composition and interchange. A strict map out is exactly a compatible pair
of strict maps from the two sides, proving the strict pushout universal
property. Smallness bounds the term construction. No embedding theorem is
part of that proof.

This does not automatically give a pushout for arbitrary pseudofunctors.
A weak syntax must explicitly include invertible unit/compositor generators,
pure-agreement cells and their equations, and must state whether its
mapping property uses icons, pseudonatural cocones or another chosen
transformation notion. Separate strictifications of the sides need not
strictify their shared diagram compatibly; see the distinction discussed in
\cite{LackCompanion}. Similarly, an evaluation from unrestricted words to
law-governed short normal forms need not be faithful.

A retraction gives a useful but weaker condition. In the strict setting, if
$R:N\to M$ satisfies $RJ_N=J_M$, the pushout universal property supplies
$r:P\to M$ with $ri_M=\id_M$. Thus $i_M$ is faithful and reflects
isomorphisms. For strict 2-categories the same argument gives local
faithfulness and reflection of invertible 2-cells, not local fullness.
A true full-recovery theorem requires every target cell between old arrows
to come from a unique old cell, as was proved directly for
probability/list. A semantic retraction alone does not establish that.

\section{What has been assembled}

The deliverable is the relative algebra $H$ with its designated pure base,
unit, multiplication, hom-category cells and recovered endpoint squares.
For context/content input, its mixed carrier and weak global comparison are
explicit; the law turns that carrier into one compositional object. The
probability/list theorem constructs the law and proves conservative
one-sided recovery in a nonthin example. Grothendieck extension and
telescopes act within, or change the base of, this same algebraic setup.
Fixed effects, Para, displays and opwreaths enter through specified
operations rather than being declared related by an unnamed common label.

The general claims have three different strengths: native operational
recovery with every square; marked recovery of supplied presentation data;
and mixed interaction conditional on a fully coherent law, with
conservativity checked separately. The bicyclic example prevents replacing
these distinctions by a universal faithful merger. Lax structural maps are
allowed in the explicitly given lax-double comparison class, not silently
substituted for pseudoinvertible algebra coherences. These are boundaries
of the mathematical construction, not missing object-level formulas.

\begin{thebibliography}{99}
\bibitem{Lawler}
F. Lawler, \emph{Fibrations of Predicates and Bicategories of Relations},
2015. Section 4.1.1, Theorem 4.1.3 and Definition 4.1.8.
\url{https://arxiv.org/abs/1502.08017}.

\bibitem{CapucciMyers}
M. Capucci and D. J. Myers, \emph{Contextads as Wreaths; Kleisli, Para, and
Span Constructions as Wreath Products}, arXiv:2410.21889v1, 2024.
Definition 3.17, Example 3.21, Theorem 4.24 and Section 6.
\url{https://arxiv.org/abs/2410.21889}.

\bibitem{Marmolejo}
F. Marmolejo, \emph{Distributive laws for pseudomonads}, Theory and
Applications of Categories 5(5), 1999, 91--147. Section 4 and Theorem 5.1.
\url{https://tac.mta.ca/tac/volumes/1999/n5/5-05abs.html}.

\bibitem{GambinoLobbia}
N. Gambino and G. Lobbia, \emph{On the formal theory of pseudomonads and
pseudodistributive laws}, Theory and Applications of Categories 37(2), 2021.
Definition 4.1 and Theorem 4.5, the latter used only as a lifting reference.
\url{https://eprints.whiterose.ac.uk/id/eprint/170302/1/pseudomonads-tac.pdf}.

\bibitem{Kamiya}
K. Kamiya, \emph{The Grothendieck construction between the telescope monad
and contextads}, Hexagon 2610.00302v1, 17 October 2026.
Construction 3.1, Proposition 3.2, Corollary 2.20 and Section 3.2.
\url{https://hexagonmath.org/2610.00302v1}.

\bibitem{Street}
R. Street, \emph{Wreaths, mixed wreaths and twisted coactions},
arXiv:1610.07281, archived version used for Definition 5 and Proposition 7.
\url{https://arxiv.org/abs/1610.07281}.

\bibitem{Garner}
R. Garner, \emph{Polycategories via pseudo-distributive laws},
arXiv:math/0606735v2. Higher two-sided Kleisli construction.
\url{https://arxiv.org/abs/math/0606735}.

\bibitem{RosebrughWood}
R. Rosebrugh and R. J. Wood, \emph{Distributive laws and factorization},
Journal of Pure and Applied Algebra 175, 2002, 327--353.
\url{https://mta.ca/~rrosebru/articles/fac.pdf}.

\bibitem{LackProps}
S. Lack, \emph{Composing PROPs}, Theory and Applications of Categories
13(9), 2004, 147--163.
\url{https://tac.mta.ca/tac/volumes/13/9/13-09.pdf}.

\bibitem{Stepan}
M. \v{S}t\v{e}p\'an, \emph{Lax structures in 2-category theory},
arXiv:2504.04467v1, 2025. Theorems 3.3.7 and 3.3.33.
\url{https://arxiv.org/abs/2504.04467}.

\bibitem{LackCompanion}
S. Lack, \emph{A 2-categories companion}, 2010.
Discussion of strict 2-functors, pseudofunctors and strict 2-category colimits.
\url{https://www.math.uchicago.edu/~may/IMA/Incoming/Lack/companion.pdf}.
\end{thebibliography}
\end{document}
